\documentclass[10pt,a4paper]{amsart}
\usepackage[margin=19mm]{geometry}
\usepackage{amsmath,amssymb,mathtools}
\usepackage[scr=rsfs,cal=euler]{mathalfa}
\usepackage{xfrac}
\usepackage[unicode,hidelinks,bookmarks=false]{hyperref}
\newcommand{\ie}{i.e.\ }
\newcommand{\cf}{cf.\ }
\newcommand{\resp}{respectively\ }
\newcommand{\Subsection}{\subsection}
\providecommand{\nobreakdash}{\nobreak\hbox{-}}
\setlength{\emergencystretch}{2em}
\multlinegap0pt
\usepackage{mdframed}
\usepackage{mdframed}
\usepackage{mdframed}
\usepackage{mdframed}
\usepackage{xcolor}
\usepackage[shortlabels]{enumitem}
\usepackage{amssymb}
\usepackage{mathtools}
\usepackage{bm}
\usepackage{dsfont}
\usepackage{xfrac}
\usepackage{tensor}
\usepackage{float}
\usepackage{algorithm}\usepackage{algpseudocode}
\usepackage{tikz}
\usepackage{tcolorbox}
\usepackage{mdframed}
\newtheorem{lemma}{Lemma}
\usepackage{cleveref}
\crefname{lemma}{Lemma}{Lemmas}
\newcommand{\EE}{\mathbb{E}}
\newcommand{\RR}{\mathbb{R}}
\title{Scaling Laws from Sequential Feature Recovery:\\ A Solvable Hierarchical Model}
\author{Arie Wortsman Zurich, Hugo Tabanelli, Yatin Dandi, Florent Krzakala, Bruno Loureiro}
\date{Source: arXiv:2605.14567v1 (CC BY 4.0)}
\begin{document}
\maketitle

\noindent\emph{Source and licence.} Scaling Laws from Sequential Feature Recovery:\\ A Solvable Hierarchical Model. Version arXiv:2605.14567v1 (CC BY 4.0), distributed by arXiv under the Creative Commons Attribution 4.0 International licence, http://creativecommons.org/licenses/by/4.0/, as recorded in the arXiv licence metadata for this exact version and shown on the abstract page https://arxiv.org/abs/2605.14567v1. Full licence notice: \texttt{licenses/arxiv-2605.14567-CC-BY-4.0.txt}.\par\smallskip

\noindent\textit{Editorial scope.} Counted unit: the article's lemma lemma:bound\_square\_y\_s (source0.tex lines 2405--2411). The article's complete original proof of it is delivered with it (source0.tex lines 2413--2461, one \texttt{proof} environment of 49 lines). No mathematics is written, repaired or re-derived: the statement and the proof are the article's own lines, in the article's own notation, and no retained line is substituted or re-flowed.  Retained uncounted as the source-local context the proof needs: lines 2026--2031.  Nothing was omitted from any retained span, so the package carries no omission record.  No retained line contains a citation, so the wrapper carries no bibliography and no bibliography entry.  No retained line contains a figure environment, an image-inclusion command or drawing code, so no image is delivered and the package declares no asset dependency.  Editorial formatting only, in addition: an AMS \texttt{amsart} preamble that restates the article's own declarations and macros used by the retained lines, namely the environments \texttt{lemma} on separate counters, together with the standard packages those lines need.  The retained environments print in this document as Lemma 1 in order of appearance, whereas the article prints the same items with the numbering of its own full document.  The complete native source of the article as distributed in the arXiv e-print for this exact version is bundled unchanged as \texttt{sources/200595/source0.tex}.
\begin{lemma}
    Let $h^{(2)} = \sum_{i=1}^{d_1} \lambda_i (I_{q}(A_i)^{2} -1)$, and let $g$ be a polynomial with information exponent $1$. Then, given $i, j \in [d_1]$ with $i \not = j$: 
    \[
     \left | \EE \left [ g(h^{(2)}) I_{q}(A_j) I_{q}(A_k)\right ] \right | = O \left ( \dfrac{Z_\gamma}{d^{q}}(j^{-\gamma} + k^{-\gamma}))\right ). 
    \]\label{lemma:non_linear_case_cross_expectations}
\end{lemma}

\begin{lemma}
    Let $k,j \in [d_1]$, with $k \not =j$. Assume $n = \omega_{d}(\frac{\min(k,j)^{2\gamma}}{Z_\gamma ^{2}}d^{q} )$. Then with high probability 
    $$
        \left ( \dfrac{Z_\gamma }{n} y_\mu h_{\mu,k} h_{\mu,j} \right )^{2} \lesssim \dfrac{Z_\gamma ^2}{d^{q}} \min(k,j)^{-2\gamma}. 
    $$
    \label{lemma:bound_square_y_s}
\end{lemma}

\begin{proof}
    We want to concentrate
    \begin{align}
         E^{2} =\left ( \dfrac{Z_\gamma }{n} y_\mu h_{\mu,k} h_{\mu,j} \right )^{2}. 
    \end{align}
     For this, we begin by decomposing: 
    \begin{align}
        E^{2} & = \left (  \dfrac{1}{n} \sum_{\mu=1}^{n} \left ( y_\mu h_{\mu,k} h_{\mu,j} -\EE\left [y_\mu h_{\mu,k} h_{\mu,j}\right ]\right ) + \EE \left [ y_\mu h_{\mu,k}h_{\mu,k} \right ] \right )^{2} \\ 
        & \lesssim \underbrace{\left ( \dfrac{1}{n} \sum_{\mu=1}^{n} \left ( y_\mu h_{\mu,k} h_{\mu,j} -\EE\left [y_\mu h_{\mu,k} h_{\mu,j}\right ]\right )\right )^{2}}_{(I)} + \underbrace{\EE \left [ y_\mu h_{\mu,k}h_{\mu,k} \right ] ^{2}}_{(II)}. 
    \end{align}
    We begin by concentrating $(I)$. For this, we note that since all terms are centered, independent random variables: 
    \begin{align}
        \EE \left [ (I) \right ] & = \dfrac{1}{n^{2}} \sum_{\mu =1 }^{n}  \EE \left [ \left ( y_\mu h_{\mu,k} h_{\mu,j} -\EE\left [y_\mu h_{\mu,k} h_{\mu,j}\right ]\right )^{2}\right ].
    \end{align}
    Note that, since $\mathrm{Var}(y_\mu) \lesssim 1$, we have: 
    \begin{align}
        \EE \left [ \left ( y_\mu h_{\mu,k} h_{\mu,j} -\EE\left [y_\mu h_{\mu,k} h_{\mu,j}\right ]\right )^{2}\right ] \lesssim \EE \left [ \left ( y_\mu h_{\mu,k} h_{\mu,j} \right )^{2} \right ] \lesssim 1. 
    \end{align}
    Then
    \begin{equation}
        \EE \left [ (I) \right ] \lesssim \dfrac{1}{n},
    \end{equation}
    and since $(I)$ is positive, by Markov inequality we conclude that with high probability: 
    \begin{equation}
        \left ( \dfrac{Z_\gamma }{n} y_\mu h_{\mu,k} h_{\mu,j} \right )^{2} \lesssim \dfrac{1}{n}  + \EE \left [ y_\mu h_{\mu,k}h_{\mu,k} \right ] ^{2}. 
        \label{eq:bound_I_sum_y}
    \end{equation}
    We now turn to $(II)= \EE \left [ y_\mu h_{\mu,k}h_{\mu,k} \right ] ^{2}$. Recall that, by definition: 
    \begin{equation}
        y_\mu = g\left ( \sum_{p=1}^{d_1}  \lambda_p ((h^{(1)}_{\mu, p})^{2} - 1) \right ),
    \end{equation}
    for some polynomial $g:\RR \to \RR$, with $\EE[g'(h^{(2)})] \not = 0$. Then, by \Cref{lemma:non_linear_case_cross_expectations}
    \begin{align}
         \left | \EE \left [ y_\mu h_{\mu,k}h_{\mu,k} \right ] \right | &= O \left ( \dfrac{Z_\gamma }{d^{\frac{q}{2}}} \min(k,j)^{-\gamma} \right ). 
    \end{align}
    Taking the square:
    \begin{equation}
        \EE \left [ y_\mu h_{\mu,k}h_{\mu,k} \right ]^{2} \lesssim \dfrac{Z_\gamma ^2}{d^{q}} \min(k,j)^{-2\gamma}. 
        \label{eq:bound_II_sum_y}
    \end{equation}
    By putting together \Cref{eq:bound_I_sum_y} and \Cref{eq:bound_II_sum_y}: 
    \begin{equation}
         \left ( \dfrac{Z_\gamma }{n} y_\mu h_{\mu,k} h_{\mu,j} \right )^{2} \lesssim \dfrac{1}{n} + \dfrac{Z_\gamma ^2}{d^{q}} \min(k,j)^{-2\gamma}.
    \end{equation}
    and since $n = \omega_{d}(\frac{\min(k,j)^{2\gamma}}{Z_\gamma ^{2}}d^{q} )$,  we conclude that with high probability: 
    \begin{equation}
        \left ( \dfrac{Z_\gamma }{n} y_\mu h_{\mu,k} h_{\mu,j} \right )^{2} \lesssim \dfrac{Z_\gamma ^2}{d^{q}} \min(k,j)^{-2\gamma}. 
    \end{equation}
\end{proof}

\end{document}
